\documentclass[tikz,border=5pt]{standalone}
\usepackage[T1]{fontenc}
\usepackage[utf8]{inputenc}
\usepackage[italian]{babel}
\usepackage{amsmath}
\usetikzlibrary{arrows.meta,calc,decorations.pathreplacing,positioning}

\begin{document}
\begin{tikzpicture}[
    % Stili principali della finestra.
    cell/.style={
        draw,
        minimum width=0.72cm,
        minimum height=0.72cm,
        font=\ttfamily\large,
        inner sep=0pt
    },
    search/.style={fill=blue!7},
    lookahead/.style={fill=green!9},
    match/.style={fill=orange!28},
    nextsymbol/.style={fill=red!18},
    pointer/.style={-Latex,thick},
    note/.style={font=\small,align=center},
    title/.style={font=\bfseries\small}
]

% Riga dei caratteri: search buffer a sinistra, look-ahead buffer a destra.
\foreach \ch [count=\i from 1] in {a,b,a,a,b,a,b,a,a,b,b} {
    \ifnum\i<6
        \node[cell,search] (c\i) at ({0.72*(\i-1)},0) {\ch};
    \else
        \node[cell,lookahead] (c\i) at ({0.72*(\i-1)},0) {\ch};
    \fi
}

% Evidenziazione del match e del simbolo successivo.
\foreach \i in {1,2,3,6,7,8} {
    \node[cell,match] at (c\i) {};
    \node[cell,match] (c\i) at (c\i) {\ifcase\i\or a\or b\or a\or\or\or a\or b\or a\fi};
}
\node[cell,nextsymbol] (c9) at (c9) {a};

% Separatore tra le due parti della finestra.
\node[title,blue!60!black] at ($(c3.north)+(0,0.45)$) {search buffer};
\node[title,green!45!black] at ($(c8.north)+(0.3,0.45)$) {look-ahead buffer};

% Graffe descrittive.
\draw[decorate,decoration={brace,mirror,amplitude=5pt}]
    ($(c1.south west)+(0,-0.12)$) -- ($(c5.south east)+(0,-0.12)$)
    node[midway,below=6pt,note] {testo gi\`a codificato};
\draw[decorate,decoration={brace,mirror,amplitude=5pt}]
    ($(c6.south west)+(0,-0.12)$) -- ($(c11.south east)+(0,-0.12)$)
    node[midway,below=6pt,note] {testo ancora da codificare};

% Puntatori i e j.
\draw[pointer,blue!70!black] ($(c6.north)+(0,1.0)$) -- ($(c6.north)+(0,0.12)$)
    node[pos=0,above,note] {$i$};
\draw[pointer,orange!75!black] ($(c1.north)+(0,1.0)$) -- ($(c1.north)+(0,0.12)$)
    node[pos=0,above,note] {$j$};

% Collegamento tra prefisso del look-ahead e occorrenza precedente.
\draw[pointer,orange!80!black,bend left=22]
    ($(c7.north)+(0,0.15)$) to
    node[midway,above=5pt,note] {match $\alpha=\texttt{aba}$}
    ($(c2.north)+(0,0.15)$);

% Spiegazione della tripla prodotta.
\node[note,align=center] at ($(c6.south)+(0,-1.45)$) {
    $o=i-j=5$ \quad $\ell=|\alpha|=3$ \quad $s=\texttt{a}$
    \quad $\Rightarrow$ \quad output: $\langle 5,3,\texttt{a}\rangle$
};
\node[note,align=center] at ($(c6.south)+(0,-2.0)$) {
    si copia \texttt{aba} dal search buffer e poi si scrive il simbolo successivo \texttt{a}
};

\end{tikzpicture}
\end{document}
